\chapter{Study cases}\label{ch:cases}

A \emph{study case} is one question asked of the network:

\begin{itemize}
  \item a \textbf{study}, such as the Thevenin equivalent or the load flow;
  \item its \textbf{options}, such as the bus to look at or the fault type;
  \item optionally, a \textbf{network variant}: elements taken \emph{out of service} and
    fields \emph{overridden} by expressions.
\end{itemize}

Cases are saved with the project, so one project holds every question about a network
side by side: ``industry disconnected'', ``load at 50\,\%'', ``+10\,\% load at bus 4''.
The variant is applied to a copy of the network when the case runs. The network you drew
never changes.

\section{The Study cases panel}

\begin{minipage}[t]{0.5\linewidth}
The panel lists the project's cases in the order \keys{F5} runs them. A case with a
variant says so after its title, for example \emph{[1 override(s)]} or \emph{[1 out of
service]}. Hover a case to see its id and study.

\begin{description}
  \item[New...] opens a menu of every study, grouped as in the \gui{Calculate} menu.
    Pick one to create a case and edit it.
  \item[Edit...] edits the selected case. Double-clicking a case does the same.
  \item[Duplicate] inserts a copy below it, with a new id and \emph{(copy)} added to its
    title.
  \item[Delete] removes it.
  \item[Run selected] runs the selected cases. Hold \keys{\ctrl} or \keys{\shift} to
    select several.
  \item[Run all (F5)] runs every case.
\end{description}
\end{minipage}\hfill
\begin{minipage}[t]{0.45\linewidth}\vspace{0pt}
  \centering\shot[0.97\linewidth]{panel-cases.png}
\end{minipage}

Every change to the list (new, edit, duplicate, delete) is undoable. A project without
cases is valid. \keys{F5} then runs every study once with its default options, which is
a quick way to see everything a new network gives.

\section{The case dialog}

The dialog shows the study's description, plus a note for studies that need one, and two
tabs.

\subsection{Options}

\begin{minipage}[t]{0.5\linewidth}
\gui{Case title} names the case in the panel and in its result tab. Left empty, the
study's title stands in.

The other rows are the study's options, generated from its definition and described for
each study in \cref{ch:studies}. Bus options list the network's buses. Element options
list the elements of the right type. Choice options offer their values. Options that only
apply in some situations appear when they do: the ground impedance of a fault, for
instance, appears for a double line-to-ground fault only.

Some options depend on others. In S1a, the list of shunts that may be included at the
ports changes with the port buses chosen (\cref{sec:study-s1a}).
\end{minipage}\hfill
\begin{minipage}[t]{0.45\linewidth}\vspace{0pt}
  \centering\shot[0.97\linewidth]{dialog-case-options.png}
\end{minipage}

\subsection{Network variant}\label{sec:variant}

\begin{minipage}[t]{0.5\linewidth}
\textbf{Out of service.} Every bus and component is listed with a check box. A checked
element is out of service in this case only. Assignment S2 uses this to study the grid
with the S1 industry disconnected (TI1 out of service in \file{grid\_industry.psn}).

\textbf{Overrides.} Each row sets a field, written \code{COMPONENT\_ID.field} (for example
\code{LD4.p\_mw}), to an expression. The expression may use the project variables and
\code{value}, the field's current value:

\begin{itemize}
  \item \code{1.1*value} raises the field by 10\,\% (Assignment LFd);
  \item \code{0.5*value} halves it (Assignment S2d);
  \item \code{A/1000} sets it from a variable.
\end{itemize}

\gui{Add override} adds a row, and \gui{Remove} deletes the selected rows. Field names
are the ones in the reference tables of \cref{ch:components}, and in the data manager's
column tooltips.
\end{minipage}\hfill
\begin{minipage}[t]{0.45\linewidth}\vspace{0pt}
  \centering\shot[0.97\linewidth]{dialog-case-variant.png}
\end{minipage}

When you press \gui{OK}, \app\ checks every override: the component must exist, the field
must be one of its fields, and the expression must evaluate. A mistake is shown in red
under the tabs and the dialog stays open. The tab's title counts the variant's entries,
for example \gui{Network variant (1)}.

\begin{note}
Overrides are applied in the order shown, on a copy of the network taken after the
project's own bindings were evaluated. When a later edit deletes or renames an element
that a case refers to, the case runs anyway and its result starts with a warning naming
the stale entry.
\end{note}

\section{Running cases}

Cases run in three ways:

\begin{itemize}
  \item \keys{F5}, \menu{Calculate > Run all cases} or \gui{Run all (F5)}: every case of
    the project, in order;
  \item \gui{Run selected}: the selected cases;
  \item a study from the \gui{Calculate} menu, which runs that study's own case
    (\cref{sec:own-case}), or a bus's context menu, which runs a Thevenin equivalent or a
    three-phase fault at that bus without saving a case.
\end{itemize}

Before running, \app\ validates the network. An error, such as a missing base reference
bus or a branch that joins a bus to itself, stops the run with a list of what to fix
(\cref{ch:troubleshooting}). A warning, such as an unconnected component, does not stop
it. Each case then gets a result tab. A case that cannot run (S1a with ports that are not
connected, a Thevenin equivalent in a network without a source) gets a tab with the
reason, and the other cases are still shown. Only when every case fails does \app\ show
the error in a message box instead.

\section{A study's own case}\label{sec:own-case}

Choosing a study from the \gui{Calculate} menu does not run it blindly. It opens the
dialog of the study's \emph{own case}, the case whose id is the study's key (\code{thevenin},
\code{loadflow}\dots), creating it the first time with the study's defaults. Accept the
dialog and the case is saved in the project and run. The next time you choose the study,
the dialog opens with the options you left. These cases appear in the \gui{Study cases}
panel like any other, so \keys{F5} runs them too. Studies without options (the per-unit
model and the sequence model) run straight away.

\section{Default options}

A study's options start from defaults chosen for the network at hand:

\begin{itemize}
  \item the bus options of the Thevenin and short-circuit studies default to the base
    reference bus, unless an ideal source sits there, in which case the bus furthest from
    it is used;
  \item S1a's ports default to the reference bus and the bus furthest from it;
  \item the linear voltages report the bus furthest from the reference bus;
  \item the unbalanced load defaults to the last load of the network, with its phase
    factors at 1, and reports the voltages at its own bus.
\end{itemize}

An option left empty in a case falls back to its default when the case runs.
